\documentclass[12pt,a4paper]{article}
\usepackage[utf8]{inputenc}
\usepackage[T1]{fontenc}
\usepackage[portuguese]{babel}
\usepackage{setspace}
\usepackage{indentfirst}
\usepackage{geometry}
\usepackage{microtype}
\geometry{top=3cm, bottom=3cm, left=3.5cm, right=3cm}
\setlength{\parindent}{2em}
\setlength{\parskip}{0.4ex}

\newcommand{\secnum}[2]{\medskip\noindent\textbf{#1. #2}}
\newcommand{\pagmark}[1]{\quad[\textit{p.\,#1}]\quad}

\begin{document}

\begin{center}
  {\large \textbf{Essência e Existência}}
\end{center}

\vspace{1.5em}

\pagmark{302}

\medskip\noindent\textbf{29. Essência e existência. — O ente, como dissemos, denota}
uma coisa, que existe ou pode existir. Por isso o ente compõe-se
de essência e de existência, como de potência e de ato. — Tendo
declarado-a..ideia..de ente e a ideia de potência e de ato, devemos
agora tratar da essência e da existência \footnote{O assumpto, de que nos ocupamos neste artigo, é da uma importância suprema em toda a filosofia. Pedimos, pois, toda a atenção dos
 nossos pacientes leitores.. 
\pagmark{302}

\medskip\noindent\textbf{30. Essência. — Essência é a coisa, pela qual um ente é O}
que é: ou, mais desenvolvidamente, é a coisa, a qual constitui um
\pagmark{303}
ente numa certa espécie, ou num certo gênero, e o distingue dos
 entes das outras espécies, ou dos outros gêneros; assim a essência
 de Pedro é a humanidade, porque é por ella que Pedro é homem,
 e não é outra coisa, — ou porque é ella  que constitui Pedro
 na espécie humana e o distingue dos animaes irracionaes \footnote{(1) São muitas e varias as explicações, que se dão da essência. Assim
 diz-se também que a essência
 — é a coisa, cujo ato é O ser, ou à existência,
 — é o que primeiramente se concebe num ente, — é o que faz conhecer
 a
 coisa como é em si mesma, — é o que distingue uma coisa de todas as ou-
 tras, — é a raiz e O sujeito de todas as propriedades e operações de um
 ente, etc. — Todavia, embora sejam muitas e varias as explicações. que se
 dão da essência, contudo todas exprimem uma coisa, concebida como o fun
 do, que constitui um ente numa certa categoria ou espécie e de que brotam,
 como do seu princípio radical, todas as propriedades e operações do
 ente.
 Apresentamos explicações da essência, e não definições. Por quanto, a
"essência, embora seja menos abstrata que O ente, contudo, pela sua, simpli-
 cidade, não é susceptivel. de uma definição. rigorosa,-seh--que-se-conimetta
uma petição de princípio. Com efeito, devendo a definição rigorosa explicar
O" que una coisa €, isto é, devendo explicar a essência de uma coisa, não é
possível formar uma definição rigorosa, sem que primeiramente se possua
 uma noção geral da própria essência. A essência, pois, é uma das noções
primitivas e immediatamente evidentes, que não precisa de demonstração,
nem pode ser demonstrada.
 A essência chama-se também quididade, forma, natureza; mas esses
vocabulos denotam sempre uma e a mesma coisa, embora sob diferentes
aspectos. — Chama-se quididade; porque a essência, constituindo as. coisas
no próprio gênero ou na. própria espécie, é denotada pela definição, com a
«qual respondemos à pergunta: o que é isto — quid est res? — Chama-se
forma; porque, assim como a forma é ato e perfeição da materia, porque
 a
tira da indiferença e a constitui num determinado grau de ser, assim à
essência é ato e perfeição do ente, que pode considerar-se como um sujeito, constituido pela própria essência num determinado grau de ser. —
 Chama-se natureza, enquanto é o primeiro princípio da operação própria, porque nenhuma coisa é destituida da sua operação. — Notamos que a essência
chama-se essência, enquanto é nella e por ella que uma coisa tem o ser
(esse). Uma coisa tem o ser (esse) ma essência; porque a essência é a potencia, em que o ser é recebido. Una coisa tem o ser (esse) pela essência; porque é por causa da essência que uma coisa tem o ser, pelo qual é constituida numa certa espécie ou gênero, e porque nenhuma coisa poderia receber
o ser, se não tivesse « capacidade ou a potência para isso,
 e a essência é
essa capacidade. — Diz S. Tomás (De Ente et Ess., c. 1): «Illud per quod
res constituitur in próprio genere vet specie, est quod significamus per
 definitionem indicantem quid est res; inde est quod nomen essentiae à philosoPhis in nomen quidditatis mulatur... Dicitur etiam forma,
 secundum quod
Per formam significatur perfectio vel certitudo
 uniuseuiusgue rei... Nomen
autem natura videtur significare essentiam rei,
 secundum quod habet ordinem vel ordinationetm ad propriam operation
 em rei, cum nulla res própria,

\pagmark{304}
destituatur operatione... Sed essentia dicitur secundum quod per cam el
 in ea res habet esse ».
 Como se vê, a ideia de essência succede immediatamente à de ente;
 porque, tendo conhecido um'objeto sob a ideia indeterminada de ente, a
 primeira pergunta, que fazemos acerca do mesmo objeto, é a seguinte: o
 que é isto? — A resposta exprime a essência.. 

\pagmark{304}

\medskip\noindent\textbf{31. Divisão da essência. — À essência é — a) possível e}
atual, — b) substancial e acidental. — c) física e metafísica.

     \medskip\noindent a) Essência possível e atual, — Possível é a que presentemente não existe, mas pode existir na ordem real. — Atual é
a que presentemente existe na realidade. — Neste artigo Lrala-
 mos da essência atual. .

     \medskip\noindent b) Essência substancial e acidental. — A substancial é a
que-se encontra na substância; tal é a essência do homem. —
 tal é à essência
 A acidental é a que se encontra no acidente
 da sabedoria, — A essência existe na substância por um modo
absoluto. — no acoidente por um modo relativo.

     \medskip\noindent c) Essência: física e metafísica. — É física, quando
. os seus elementos são concretos, reais e realmente distintos; assim
 do homem é constituida pela alma racional
 a essência física
 e pelo corpo organico. — É metafísica, quando os seus ele-
 mentos-são-abstratos; togicos-e-logicamente distintos, denotando
 o gênero e a diferença; assim a essência metafísica do ho-
 mem é constituida pela animalidade e pela racionalidade \footnote{A essência física é artificial ou natural. É artificial, se a sua con-
 stituição depende do artista; tal é a essência de uma relogio. É natural,
 quando a sua constituição é obra da natureza; tal é a essência de uma
 planta. — A essência metafísica é apenas o conceito abstrato da própria
 essência física, dividido em dois oufros conceitos, que são O gênero e x
 diferença. Como dissemos na Lógica. O gênero não exprime uma parte, física
 ou realmente distinta das outras, mas exprime a essência completa, ainda
 que de um modo indeterminado. Assim, na definição do homem, O gênero ani-
 mal não denota exclusivamente o corpo; mas denota o corpo e à alma, em-
 quanto à perfeição e ás forças da vida vegetativa e sensitiva, abstrahindo da
 perfeição e das forças da vida intellectual, contidas no conceito de racional,
 que exerce o munus de dijferença (Cf. De Ente et Ess. c. VI). — Mas as partes,
 quer sejam físicas, quer sejam »metafísicas, sempre se referem à essência
 real; pois a essência física e é metafísica são uma e a mesma essência
 real, considerada sob um duplice aspecto — nas coisas e na inteligencia.. 
\pagmark{304}

\medskip\noindent\textbf{32. A essência é uma coisa real e objetiva.}

     \medskip\noindent a) A essência É o que constitui o ente na, própria espécie,
 ou no próprio gênero, e o distingue dos entes das outras espe-
 cies, ou dos outros. gêneros. Ora o que constitui o ente na
\pagmark{305}
 própria espécie, ou no próprio gênero, e o distingue dos entes
das outras espécies, ou dos outros gêneros, não 'é um mero
producto da nossa fantasia, mas é uma coisa real e objectiva;
porque as coisas são o que são, independentemente dos nossos
atos. Logo a essência é uma coisa real e objectiva.

     \medskip\noindent b) Toda a essência é física ou metafísica. Se é física,
é certamente real e objectiva; assim o corpo organico e a alma
racional são coisas reais e objetivas, e por isso real e objectiva
deve ser a essência, que é constituída por esses dois elementos.
— Se é metafísica, é também real e objectiva; porque a essência
metafísica, sendo a noção abstrata da essência física, é
uma coisa que — ou existe na realidade, embora por um modo
differente do modo, por que existe na inteligência, — ou na realidade tem o seu fundamento, e por isso é sempre uma coisa
real e objetiva. Logo a essência é uma coisa real e objectiva.\footnote{Veja-se o que dissemos na Lógica com relação à objectividade das
ideias universais. — São Os sensualistas e os positivistas, que reduzem a
essência das coisas à mera ficção da nossa fantasia, a nomes ou conceitos
exclusivamente subjectivos, aos quaes nenhuma realidade corresponde.}

\medskip\noindent\textbf{33. A essência encontra-se não só na substância, mas tam-}
bem no acidente. — A essência é a coisa, pela qual o ente é
o. que é, — ou que constitui o ente numa certa espécie, ou num
certo gênero, e o distingue dos entes das outras espécies, ou
dos outros gêneros. Ora não só a substância, mas também o accidente é o que é, pertence a uma certa espécie, ou à um certo
gênero, e distingue-se dos entes das outras espécies, ou dos
outros gêneros. Logo a essência encontra-se não só na substância,
mas também no acidente \footnote{A conclusão é certa. A essência responde à pergunta: o que é isto?
Ora esta pergunta pode fazer-se não só com relação à substância, mas também com relação ao acidente. Por quanto, não só a substância,
 p. ex., o homem, é o que é, e distingue-se das plantas, das pedras, etc.: mas também
o acidente, p.ex., a alvura, é o que é, e distingue-se das outras córes, da
sabedoria, do poder, etc. — Todavia, a essência compete primeiramente à
substância, e secundariamente ao acidente. Na verdade, a essência deriva do
ser (esse), que existe nella e por ella. Ora o ser primeiramente compete à
substância, e secundariamente ao acidente; visto que o ser do acidente
supõe o ser da substância, e neste se funda, isto é, o ser acidental supõe
o ser substancial e neste se apoia. Por isso, à essência substancial é uma
forma, pela qual um ente é constituido no seu ser primeiro, que é o substancial,
e existe independentemente de toda e qualquer forma anterior, tal é a humanidade; 20 passo que a essência acidental é uma forma. pela qual um
\pagmark{306}
 ente, já constituido no seu primeiro ser, é diversamente modificado no seu
 ser secundario, que é o acidental, tal é a beleza, a sabedoria, a bondade, etc.
 — Esta é a doutrina de S. Tomás. Diz o Santo que a essência se deve
 atribuir à substância e ao acidente, como se atribui o ente. Ora o ente at-
 tribui-se primeira e absolutamente à substância, O secundaria e relativamente
 ao acidente. Logo também « essência deve atribuir-se primeira e absoluta-
 mente à substância, e secundaria é relativamente ao acidente. (Clr. De Ente
 et Essentia, c. III).. 
\pagmark{306}

\medskip\noindent\textbf{34. A essência substancial não é uma simples coleção}
"de qualidades.

     \medskip\noindent a) A essência substancial denota a realidade, que constitui
 o ente no seu ser primeiro e o faz subsistir em si mesmo; ao
 passo que a qualidade, como todo e qualquer acidente, denota

 uma coisa, que não constitui o ente no seu ser primeiro, mas
 sobrevem ao ente  já constituido e subsistente em si mesmo. Ora
 o que se une a um ente, já constituido no ser primeiro, não pode
 entrar, como elemento, nessa primeira constituição. Logo a es-
 sencia substancial não é uma simples coleção de qualidades.

     \medskip\noindent b) Se a essência substancial fosse apenas uma coleção de
 qualidades, quando a substância se muda nas suas qualidades,
"deveria mudar-se também na sua essência. Ora isto não sue-
 cede; porque, enquanto as qualidades mudam, apparecem e
 desaparecem, a essência permanece sempre a mesma; assim 0
 homem de bom torna-se mau, mas é sempre homem. Logo a
 essência substancial não é uma simples coleção de qualidades \footnote{A verdade da proposição resulta da própria inspeção dos termos. A
 substância subsiste em si e por isso, para existir, não precisa de sujeito, a
 que adhira; ao passo que a qualidade,
 por ser acidente, não subsiste em
 si e por isso, para existir, precisa de um sujeito, a que adhira. Não pode,
 pois, afirmar-se, sem contradição, que a essência substancial é apenas uma
coleção de qualidades. — Nem se diga que, se uma qualidade não pode
 subsistir em si mesma, pode subsistir a coleção das qualidades; porque a
 não-subsistencia é o caracter essencial das qualidades, e por isso, se a subsi-
 stencia não se encontra numa qualidade, nem se encontra na coleção de
 qualidades. Se um cego não vê, nem podem ver mil cegos.
 Foi Locke, seguido pelos, materialistas, quem ensinou — que a substan-
"cia, ou essência substancial, não é uma coisa distinta das qualidades sen-
 siveis, — e que, por isso, uma substância, concebida como um princípio ou
 sujeito, de que resultam ou a que adherem as qualidades, é uma ficção
 gratuita da nossa inteligência. — O erro de Locke é uma consequencia do
 seu systema, que admitte a sensação como única fonte de todos os conhecimentos humanos. — J falsa a afirmação de Locke, dizendo que, quando
 concebemos a essência substancial, concebemol-a como um cumulo ou uma
 coleção de qualidades; porque à nossa inteligência, quando, advertida
\pagmark{307}
 pelos sentidos, percebe as qualidades, não pára ahi, mas, através das qualidades e por meio d'elas, chega até ao fundo do objeto, e descobre a
essência, de que resultam e a que adherem as qualidades.. 
\pagmark{307}

\medskip\noindent\textbf{35. Nem todas as essências nos são desconhecidas. — A}
experiencia atesta que possuimos uma ideia clara e distinta
de muitos seres; porque não só não confundimos um com o
outro, mas também indicamos os elementos, que os constituem
num certa e determinada espécie. Ora não poderíamos extremar
um ser de outro, nem indicar os seus elementos constitutivos,
se a essência à eles nos fosse desconhecida; pois é pela essência
que um ser é o que é, e se distingue dos outros seres. Logo
nem todas as essências nos são desconhecidas \footnote{Dizemos que nem todas as essências nos são desconhecidas, porque
concedemos que muitas escapam ao nosso alcance. Acrescentamos que
mesmo a noção, que temos de algumas d'elas, embora seja clara e distinta,
é contudo inadequada e imperfeita; visto que a essência, sendo o que de
mais intimo se encontra nos entes, escapa a uma intuição directa e immediata, e só é percebida pelo raciocinio, que, baseando-se no princípio de
causalidade, sobe das propriedades e dos fenômenos, que nos são manifestos, para o princípio, de que derivam e que é a essência. — Mas, ainda que
imperfeito, o conhecimento da essência é científico; porque O método é
científico, quando se levanta, pela indução, dos efeitos para a causa, das
acções para a natureza do agente, dos fenômenos para à essência.
 . Locke, Bonnet, Condillac dizem que o homem, se pode conhecer as essencias nominais, não pode conhecer as essências reais. — Essências reais são
as que física e intrinsecamente constituem os entes, e por isso existem
independentemente do nosso pensamento; enquanto as essências nominais
são as que a inteligência compõe por meio das ideias de gênero e de diferença específica, e por isso são meros productos da imaginação, destituidos
de todo o fundamento real.
 Esta opinião é falsa, não só porque sustenta que não podemos conhecer,
de modo algum, nenhuma essência real, mas também porque afirma que as
essências nominais, por serem compostas de gênero € de diferença específica,
são meros productos da nossa imaginação, sem fundamento real. O gênero
ea diferença não são uma ficção inutil da nossa inteligencia, mas são
noções universais e abstratas, representativas de uma coisa, que é real,
ou que tem fundamento na realidade, como suficientemente explicamos,
\pagmark{308}. 
\pagmark{307}

\medskip\noindent\textbf{36. A essência das substâncias corporeas é composta. —}
As substâncias corporeas, de que é constituido o mundo, convêm entre si enquanto são substâncias, ma diferem enquanto
uma substância pertence a uma espécie, e outra a outra espécie.
Ora esta semelhança e esta diferença demonstram claramente
que a essência "dellas é composta de dois elementos: um commum, de que deriva a semelhança e que é O gênero, outro pro-
\pagmark{308}
prio, de que deriva a diferença e que é a diferença específica.
 Logo a essência das substâncias corporeas é composta \footnote{As substâncias criadas são corporeas e incorporeas. A essência das
 substâncias corporeas, — enquanto existe na realidade, é composta de dois
 princípios, ou
 de duas causas intrínsecas: uma material é outra formal, que
 constituem um composto real; — e, enquanto é conhecida e definida pela
 nossa inteligência, resulta de dois conceitos, um dos quaes é conceito
 do gênero e outro é o conceito de diferença específica, que constituem um
 composto intencional. O gênero funda-se no princípio material; a diferença,
 no princípio formal. — A essência das substâncias incorpóreas, — enquanto
 existe fora da nossa inteligência, é uma forma inteiramente simples, — mas,
 enquanto é conhecida e definida por nós, está sujeita à uma composição.
 Por quanto, a nossa inteligência, pela sua imperfeição natural, apreende
 as coisas simples e espirituais pelo mesmo modo, por que apreende as
 essências malerines e compostas, e por isso define as formas subsistentes
 pelo gênero e pela diferença, como define us essências materizes. O gênero
 funda-se na potência, a diferença funda-se no ato; porque as substâncias
 simples, embora não sejam compostas quanto à essência, contudo são real-
 mente compostas de potência e de ato, isto é, da essência e de existência.. 
\pagmark{308}

\medskip\noindent\textbf{37. Dotes das essências criadas. — As essências criadas}
 são necessárias, indivisíveis e eternas. *

     \medskip\noindent a) São necessárias, ou imutáveis, — não com relação á
 existência (porque, neste sentido, são contingentes e mutaveis), —
 mas com relação aos elementos especificos, que as constituem
 ;
 porque uma essência não pode ser nem entender-se sem os ele-
 mentos, de que é composta. Assim a essência do homem é ne-
 cessaria, ou imutável, com relação à animalidade e á raciona-
 lidade; pois não pode conceber-se um homem que não seja
 animal e racional. Logo as essências criadas são necessárias \footnote{Toda à essência pode considerar-se — ou nos seus elementos espe-
 cificos e abstratos,
 — ou nos seus elementos reais e concretos, conforme é
 metafísica ou física. — Se a essência se considera nos seus elementos
 especificos, isto é, na sua realidade quidditativa, abstrata das condições,
 que a individualisam ou determinam neste ou naquele indivíduo, importa
 no seu conceito uma coisa absolutamente necessária e imutável: porque
 reproduz em si mesma o exemplar, que existe na inteligência de Deus e que
 é necessario e imutável..  — Se, porém. a essência se considera hos ele-
 mentos reais, de que é composta nos individuos singulares, nesse sentido
 não é necessária nem imutável, porque pode ser alterada é corrompida pela
 acção dos agentes naturaes. Assim a essência real ou física do homem
 acaba com a morte, quando a alma se separa do corpo. — Mas, quando a
 nossa inteligência considera a essência abstrata dos princípios individuacs,
 percebe unicamente a realidade quidditativa, os elementos especi
 Como se vê, a necessidade, que atribuímos à essência metafísica das
 criaturas, — não é absoluta, porque Deus podia formar outros modelos ou
\pagmark{309}
exemplares, que fossem reproduzidos na criação, — mas é hipotética, em-
 quanto, suposta a criação, Deus não muda o seu plano, mas tudo executa,
 em conformidade com a sua infinita Sabedoria.. 
\pagmark{309}

     \medskip\noindent b) São indivisíveis, enquanto não podem receber ou perder
 um elemento, sem que fiquem destruidas. Tal indivisibilidade
- das essências é uma consequencia da sua necessidade. De facto,
 se a uma essência se acrescentar ou tirar um elemento, já não
 permanece a mesma, deixa de ser o que é; assim, se á substância
 viva e sensitiva, que é o bruto, se acrescenta racional, não temos
 o bruto, mas homem, e se á mesma substância se tira o ele-
 mento sensitivo, não temos o bruto, mas a planta. Logo as essen-
 cias criadas, embora compostas, são indivisíveis \footnote{As essências criadas, tanto físicas como metafísicas, por serem
 compostas, são divisíveis nas partes respectivas. As essências físicas das
 substâncias corporeas são divisíveis na materia e na forma substancial; e
 as essen: metafísicas são divisíveis, pela análise da razão, nas partes
 da definição, que são o gênero e à diferença específica. Não é, pois, neste
 sentido que as essências criadas são indivisíveis. Portanto, as essências
 podem dizer-se indivisíveis, enquanto não admitem o mais nem o menos,
 enquanto os seus elementos não são susceptiveis de adição ou de dimi-
 nuição, sem que uma essência deixe de ser o que é. Por isso as essências
 foram comparados aos números; porque, assim como um número muda de
 espécie pela adição ou subtração de uma unidade, assim também as
 essências mudam de espécie pela adição ou subtração de um elemento.
(Sum. Th. 1-2, q. 52, a. 1).. 
\pagmark{309}

     \medskip\noindent c) São eternas, — não enquanto existem realmente desde
 a eternidade (pois que, neste sentido, só Deus é eterno), — mas
 enquanto, consideradas em si e em abstrato, prescindem de
 toda e qualquer diferença de tempo, e também enquanto exi-
 stem, desde a eternidade, nas ideias arquétipos do Criador. As-
 sim foi, é e será sempre verdade que um triângulo tem tres
 angulos. Logo as-essências-criadas-são-elernas-(1): \footnote{As essências criadas podêm considerar-se em si mesmas e na Jhtel-
 tigencia divina. — Consideradas em si mesmas, as essências prescindem de
 toda a limitação do tempo e do lugar. É uma eternidade negativa. — Con-
 sideradas na inteligência divina, as essências criadas são dotadas de todos
 os atributos, próprios da Divindade: porque, nesta hipótese, as essências
criadas não se distinguem realmente da Inteligência divina, nem, por isso,
 da Essência divina, mas são a própria Essência infinita, enquanto pode ser
 diversemente imitada e participada pelas criaturas.. As essências criadas,
 pois, são outras tantas participações. ainda que inadequadas,
 da única e
simplíssima Essência infinita. Esta participação não se faz por emanação,
 mas por semelhança, enquanto o efeito é semelhante à sua causa}
\pagmark{309}

\medskip\noindent\textbf{38. Existência. — Existência é O ato, ou a realização da}
\pagmark{310}
essentia; ou, mais claro, é a perfeição, que tira q essência, em
que é recebida, do estado de possibilidade, colocando-a no estado
de atualidade. — Por isso, o ente real ou atual é um composto de essência e de existência \footnote{A existência, como a essência, é uma noção simplíssima e por si
evidente, e, à semelhança das outras noções simples e evidentes, não é capaz nem precisa de uma definição rigorosa. Por isso a nossa não é senão
uma explicação. — A existência é e diz-se ato da essência. O ato, que, —
no sentido vulgar. significa operação, — no sentido rigoroso, significa uma
coisa, que determina ou aperfeiçoa. Ora a essência é por si indiferente
quanto à existência, pode recebê-la, e por isso é um sujeito indeterminado
e determinável, imperfeito e perfectível; e é a existência, que, recebida nella,
tira-a da indiferença ou indeterminação, dando-lhe o complemento, que lhe
falta na linha de ente ou da atualidade, e assim faz com que à essência
subsista em si mesma, na realidade, e de possível se torne existente. Antes
de receber à existência, a essência era possível; recebida a existência, à esssência é uma coisa real ou atual, e opõe-se à essência possível.
 A existência;
 ou"o-ato de ser éo-ato; ou-a atualidade de todos os
outros atos, e por isso é a perfeição de todas as perfeições. E nada mais
exacto. Por quanto, todo o ato, e por isso toda a perfeição (que é ato), se
é ato considerado em si mesmo, é contudo potência considerado em relação. -
ao ato de ser; porque: nos entes finitos, nenhum ato, e por isso nenhuma
perfeição encerra na sua essentia a existência, ou O ser, isto é, o ato pode
existir, mas não existe essencialmente. E, como o ato denota atualidade e
perfeição, podemos e devemos dizer que o ato de ser. ou a existência. é a
atualidade de todos. os-ontros atos, e por isso é a perfeição de todas as
perfeições. — Ao ser nada pode acrescentar-se que seja mais formal, mais
atual, e que o aperfeiçoe e determine, como o ato determina a potência...
Por isso, o ser não é determinado por outra coisa, como a potência é determinada pelo ato; mas só pode ser e é determinado pela potência. E assim
um ser distingue-se de outro ser, enquanto um é recebido numa essência
ou natureza, e Outro noutra (o ato é determinado ou limitado pela potência,
em que é recebido). Eis as palavras do S, Doutor: «Hoc quod dico esse est
actualitas omnium: actuum, et propter hoc est perfectio omnium perfectionum. Nec intelligendum est, quod ei quod dico esse aliquid addatur, quod
sit eo formalius, ipsum determinans, sicut actus potentiam. Unde non sic
determinatur esse per aliud, sicut potentia per actum, sed magis sicut actus
per potentiam... Et per hunc modum hoc esse ab íllo esse distinguitur, in
"quantum est talis vel talis naturae» (De potentia, q. 7, a. 8 ad 9; Cfr.
Sum. Th., p. 1. q.4, a. 1.ad 3). — Sendo a existência o que de mais atual
se encontra nas coisas, rectamente os escholasticos definiram a existência:
a última atualidade de uma coisa; porque uma coisa recebe a última atualidade, isto é a última perfeição, o último complemento. quando recebe.a.
existência. — Por isso, se a essência é boa, é boa não por si, mas pelo
ser. O ser é o fundamento e a razão de todo o bem, ou de toda a perfeição,
que se encontra nas coisas. Diz 8. Tomás: « Omnis nobilitas cuinscumque
rei est sibi secundum suum esse; nulla enim nobilitas ex sua sapientia (id
est, ex ipsa essentia sapientiae), nisi per eam sapiens esset, et sic de aliis.
\pagmark{311}
Sic ergo secundum modum, quo res habent esse, est suus modus in mobilitate» (C. Gent., 1, 98; cf. De Ver, q. 21,a. Lad 4;a. ad 5; Sum.
Th., p. 1, q. 62, a. 3 ad 2 et 3).
 Como se vê, a ideia de existência succede immediatamente à de essência;
porque, depois de termos concebido o que uma coisa é (quid est), logo em seguida perguntamos se ella existe (ax est). — Todavia, do facto de a existência
ser uma coisa distinta da essência, não se segue que à existência e a essência
sejam dois entes. A essência, distinta da existência, não é; nem é a existência,
distinta da essência; mas é O ente composto de essência e de existência. —
Mas não devemos imaginar que o composto de essência e de existência seja.
uma terceira coisa, resultante d' esses dois elementos e diversa de cada um
deles, como o homem, composto de alma e de corpo, é uma terceira coisa,
distinta de cada um desses dois elementos. Quando dizemos que a essência
é, o verbo é denota a atualidade da essência, e a essência é o sujeito dessa
atualidade, e assim temos a essência existente, que não é uma terceira coisa,
distinta da existência e da essência, mas é à mesma. essência enquanto
actuada pelo ser ou pela existência (S. Tomás., Quodl. III, a. 3 ad 1).
 Notamos que o ser e a existência, embora se empreguem como termos
equivalentes, contudo não significam uma e a mesma coisa. Porquanto, o
ser significa tudo o que um ente é e possue, modelado e vasado no molde
da essência, e por isso significa toda a perfeição, de que o ente é dotado,
incluida a existência; ao passo que à existência por si indica apenas a posição do ente fora das suas causas, e vão tudo o que o ente possue.}
\pagmark{311}

\medskip\noindent\textbf{39. A existência encontra-se na substância e no acidente.}
— A existência é a última atualidade de todo o ente real. Ora
todo o ente real, que existe no mundo, é substância ou acidente.
Logo a existência encontra-se na substância e no acidente. —
Por isso, assim como ha uma essência substancial e uma essência
acidental, assim também ha uma existência substancial e uma
existência acidental. — Todavia, a existência principalmente
convem á substância, e secundariamente ao acidente \footnote{A conclusão é evidente. No mesmo individuo humano, uma coisa
é a existência ou O ser de homem, e outra coisa é a existência ou o ser de
sabio, de bom, etc. — A existência, como a essência, primeiramente compete
à substância e secundariamente ao acidente; porque o acidente tem o seu
ser ou a sua existência no ser ou na existência da substância. — Em cada
substância, não pode haver senão um único ser substancial, porque não pode
haver nella senão uma única forma substancial, da qual deriva O ser substancial, mas pode haver varios seres acidentais, porque pode haver nella
varias formas acidentais, de cada uma das quaes deriva um ser arcidental.
Assim Pedro, — que, em virtude da sua única forma substancial, isto é,
da única alma racional, possue o ser substancial, pelo qual subsiste enquanto
homem, — em virtude de varias formas acidentais, possue varios seres
acidentais, pelos quaes é branco, é virtuoso, é sabio, ete. (Cfr. Cajet., Com.
in Op. De Ente et Hss. c. 7, q. 16). — Advertimos que, nos entes compostos
de materia e forma, enquanto O ser, ou à existência, se diz derivar da forma
para o mesmo composto, enquanto a forma, aperfeiçoando e completando a
\pagmark{312}
essência material na linha da essência, dá à mesma essência à ullima disposição, necessária e suficiente para a recepção da existência, isto é, torna
a mesma essência proximamente capaz de receber a existência. Isto vale,
proporcionalmente, tanto para a forma substancial, quanto para a forma
acidental. Da forma substancial, que actua a materia prima, deriva para
 o sujeito o ser substancial; da forma acidental, que actua o sujeito já
constituido no seu ser substancial, que é o primeiro ser, deriva para elle o
ser acidental, que é o segundo ser.. 
\pagmark{312}

\medskip\noindent\textbf{40. Nos entes finitos, a existência refere-se à essência,}
como o ato à potência. — Quando ha duas coisas, uma das
 quaes aperfeiçõa e completa a outra, a primeira refere-se à segunda, como o ato se refere á potência; porque ato é e diz-se
 a coisa, que aperfeiçõa e completa, e potência é e diz-se a coisa,
 que recebe a perfeição e o complemento. Ora, nos entes finitos,
 a existência aperíeiçõa e completa a essência, porque a faz sub-
 sistir na atualidade. Logo, nos entes finitos, a existência. refe-
 re-se á essência, como O ato á potência \footnote{(1) É o argumento de S. Tomás, que estabelece o seguinte princípio:
« In 'quocumque inveniuntur aliqua duo quoram unum est complementum
alterius, proportio unius ad alterum est sicut proporlio potentiae ad actum;
nibil enim completur nísi per proprium attum» (C. Gent., J. II, e. 53). —
E noutro lugar: « Quiddilas est sicut potentia, et suum esse acquisitum est
sicul actus
 » (lo IT Sent., dist. 3. q. 1, a. 1). — Sendo a existência para à
essência O que O ato é para à potência, tudo o que dissemos com ralação
ao ato e à potência devé applicar-se á existência e à essência. — D' onde
se segue que a existência é simplesmente mais perfeila do que a essência.
"Com efícito, a existência, sendo a última atualidade da essência, enquanto
a faz passar do estado de ente em potência para o estado de ente em. ato,
é para ella um princípio de complemento e de perfeição. Ora o ente, que
aperfeiçoa é simplesmente mais perfeito do que o ente, que é aperfeicoado.
 Dissemos : simplesmente, porque, sob um certo aspecto, a essência é mais
perfeita que a existência. De facto, a essência, por limitar e determinar a
existência, refere-se à própria existência, como o que limita e determina se
refere ao que é limitado e determinado. Ora
 o determinante é mais perfeito
que o determinável. — Alem disso, O ser, ou a existência, embora não possa"
dizer-se simplesmente aperfeiçoada pela essência (a essência é aperfeiçoada
pela existência, da qual recebe a última atualidade), contudo, quando é
recebida na essência, é por esta de algum modo aperfeiçoada, enquanto
 se
diz que a existência, ou que o ser é mais ou menos nobre e perfeito, conforme
é mais ou menos nobre e perfeita a essência, em gue é recebido; assim o
ser da planta é mais perfeito que o da pedra, o ser do homem é mais perfeito que o do animal, — A razão disso é porque o primeiro Agente produz
o ser, não como subsistente em si mesmo, mas em ordem a uma certa
essência, proporcionado a uma essência (pois deve haver proporção entre a
coisa recebida e o recipiente): e por isso o ser, embora receba de Deus,
 e
não da essência, a própria perfeição, contudo recebe-s de Deus em
 ordem à
\pagmark{313}
essência; e Deus produz um ser mais ou menos perfeito, conforme a quali-
 dade da essência, na--qual-o ser é recebido -como-numa-—forma-ou-molde
 Assim, a nobreza e a perfeição «a essência é o criterio. o sigoal da nobreza
 e da perfeição do ser.. 
\pagmark{313}

\medskip\noindent\textbf{41. A existência reduz-se à categoria da essência. — Tudo}
 o que se referc a uma essência, como ato à potência, reduz-se
 á categoria, em que se encontra a própria essência; porque o
 ente coloca-se na categoria, em virtude da sua essência, e a
 potência e O ato pertencem ao mesmo gênero, ou á mesma ca-

 tegoria. Ora a existência refere-se á respectiva essência, como
 ato à potência; porque a existência, como dissemos, é a última
 atualidade, que completa a essência e a coloca na ordem real.
 Logo a existência reduz-se à categoria da essência. — Por isso,
 a existência da substância reduz-se á categoria da mesma sub-
 stancia, e é uma coisa substancial; e a existência do acidente
 reduz-se.à categoria, em que se encontra o mesmo acidente, e
 é uma coisa acidental \footnote{(1) Referindo-se a existência à essência, como o ato à potência, poderia.
 perguntar-se: A existência será apenas um acidente da. substância, ou seria
 uma parte substancial? — A esta pergunta damos a seguinte resposta, indicada
 no texto. A existência não é nem pode dizer-se acidente da substância, se
 O acidente se toma no sentido rigoroso, enquanto significa uma coisa. a
 qual, para existir, precisa de um sujeito. a que adhira. Na verdade, a exi-
 stencia dá a última atualidade à substância. colocando-a na ordem real.
 Ora é absolutamente impossível que a substância receba do acidente a última
 atualidade, — Mas, se 0 acidente não se toma em sentido rigoroso, se por
 acidente se entende tudo o que não constitui a essência de uma coisa,
 neste sentido pode dizer-se que a existência é acidente da substância. Por
 quanto, a existência não entra na essência de uma coisa, não é um elemento
 constitutivo da essência (finita), mas é a perfeição que dá a última actuali-
 dade à essência, já  constituida e completa na própria linha ou ordem da
 essência, fazendo: com que a mesma essência exista na ordem real. Portanto
 devemos dizer que a existência é parte substancial da substância, ou é ato
 substancial. — não enquanto é um elemento constitutivo da própria sud-
 stancia, ou da essência substancial, — mas enquanto dá à substância a última
 atualidade, o último complemento na linha da realidade. Por isso, a exi
 stencia não se reduz à categoria de substância, como a forma substancial
 se reduz à categoria do composto, — mas como o ato se reduz à categoria
 da potência (CI. De Pot., q. 5, a. 4 ad 3)
 Agora uma advertencia muito importante. Dizemos que à existência,
 sendo o complemento. a última atualidade da substrncia e do acidente,
 se reduz ao gênero de substância ou de acidente. Todavia, propriamente
 fallando, a existência não é nem substância nem acidente categorico; mas
 é uma coisa, que transcende toda à substância e Lodo à acidente, porque é
 ato do ente, que não é gênero e transcende todas as categorias.. 
\pagmark{314}

\medskip\noindent\textbf{42. O ente finito é composto de essência real e de exi-}
 stência real. .

     \medskip\noindent a) À essência do ente finito, considerada em si mesma, 
 necessária e eterna, e, Todavia, considerada no mesmo ente,
 contingente e temporânea. Ora a essência do ente finito não é
 contingente e temporânea senão pela existência. Logo o ente finito
 é composto de essência real e de existência real.

     \medskip\noindent b) O ser (ou a existência), considerado em si mesmo, é uno,
 indiviso, ilimitado, não-causado, e todavia, considerado nos entes
 finitos, é múltiplo, díviso, limitado, causado. Ora essa modifi-
 cação, que o ser recebe nos entes finitos, só pode derivar da
 essência, a qual, sendo por si determinada, é princípio de deter-
 minação, e por isso de divisão, de multiplicidade, de limitação
 e de dependencia. Logo o ente finito é composto de existência
 real e de essência real \footnote{(1) O ser é, por si, infinito e ilimitado; porque à limitação é O mão-ser.
 Portanto, a limitação, de que o ser se acha. cercado nos entes finitos, deriva
 da essência, a qual, por ser determinada e limitada em si mesma, determina
 e limita o ser, que ella recebe. Assim, quando digo: o ser do homem. não fallo
 do ser em si, do ser puro, que é só ser, mas fallo do ser próprio deste ente,
 e porisso de um ser limitado e determinado. — Como também, o ser é,
 por si, uno e indiviso. Se, pois, se acha diviso e multiplicado nas criaturas,
 essa divisão e multiplicação deriva de outro princípio, que é a essência (GC.
 Gent.. E, 52). Por quanto a multiplicação supõe a divisão, e a divisão
 importa uma opposição. Ora tal opposição não pode dar-se da. parte do pro-
 prio ser, mas só da parte das essências, que, sendo por si determinadas, são
 distintas, opostas, múltiplos, conforme imitam, mais ou menos, à infi-
 nita e perfeitíssima Essência divina. Diz S. Tomás: « Determinatio est plu-
 ralitatis principium primum» (Quodl. VE, a 6), — Finalmente, o ser das
 criaturas deriva de uma causa extrinseca. Ora o ser não é produzido por si,
 enquanto é ser. Diz S, Tomás: « Esse ab alio causatum non competit enti
 inquantum est ens; alias omne ens esset ab alio causatum, et sic oporteret
 procedere in infinitam in causis, quod est impossibile » (1. c.). Logo, se o
 ser des criaturas não é produzido por si, é produzido por causa da essência.
 em que é recebido. -. Logo, o ente finito é composto de um duplice princípio
 real, ou atual, que é à essência e a existência (ou O ser). .
 Quando se diz que o ente finito é composto de essência real e de exi-
 stencia ou de ser real, não deve entender-se que a realidade da essência,
 seja uma coisa distinta da realidade do ser, de modo que o ente resulte de
 duas realidades. No ente tinito, o composto de essência e de ser ha uma só
 realidade. Esta única realidade é própria do ser, pertence ao ser, e pelo ser
 é comunicada à essência; por forma que o ser É real por si, é a essência
 se torna e é real pela realidade do ser, assim como a potência se torna e é
 atual pelo ato, enquanto recebe em si O ato.. 

\pagmark{315}

\medskip\noindent\textbf{43. No ente finito, a existência distingue-se logicamente}
da essência. — Uma coisa distingue-se logicamente de outra,
quando o conceito da primeira é diverso do conceito da segunda.
Ora o conceito da existência é diverso do conceito da essência;
porque — a essência significa a coisa, que constitui o ente numa
determinada espécie, ao passo que a existência significa a coisa,
que faz existir a mesma essência na realidade, — a essência é
capaz de existir, mas é por si indeterminada quanto á sua
presença na ordem real, e por isso tem razão "de potência, ao
passo que a existência tira da indeterminação a essência, preenche aquella capacidade, e por isso tem razão de ato. Logo, no
ente finito, a existência distingue-se logicamente da essência \footnote{(1) Que, em todo o ente finito, a existência seja logicamente distincla
da essência, é verdade admittida por todos os que sabem o que é distineção
lógica e o que significam os termos — existência € essência. — Como se vé,
falamos da essência real, ou atual, de que todas as criaturas existentes
são intrinsecamente compostas, e não da essência possível, que pode existir
mas que atualmente não existe na realidade.. 
\pagmark{315}

\medskip\noindent\textbf{44. No ente finito, a existência distingue-se realmente da}
essência.

     \medskip\noindent a) Como o nosso conhecimento é uma reprodução ideal
das coisas, se, no ente finito, a existência é à essência não sé
distinguissem realmente, mas fossem uma e a mesma coisa, a
espécie ou imagem inteligível, que nos representa a essência da
coisa entendida, deveria representar-nos também a sua existência;
porque uma coisa não pode separar-se de si mesma. Ora a experiencia atesta que uma espécie ou imagem inteligível pode
representar e efectivamente representa a essência das coisas,
prescindindo contudo da sua existência. Logo, no ente finito, a
existência distingue-sê realmente da essência (2).\footnote{É evidente, Se a existência fosse um elemento constitutivo da essencia, toda a definição da essência de um ente deveria exprimir também a
sua existência; visto que a definição não pode deixar de exprimir todos os
elementos, que constituem a essência do ente definido. Ora a definição da
essência de um ente prescinde completamente da existência ou da não-existencia do mesmo ente; porque podemos entender o que é o homem, e todavia ignorar se existe na realidade. Logo a existência, no ente finito, não
é um elemento constitutivo da essência, mas uma distingue-se realmente de
outra, — Por outras palavras, se o ser da criatura está fora da essência
quanto à inteligência, deve estar também fora da essência quanto à realidade.
— Diz S. Tomás: «Quidquid non est de intellectu essentiae vel quidditatis, hoc est adveniens extra et faciens compositionem cum essentia; quia
\pagmark{316}
 nulla essentia sine his, quae sunt partes essentiae, intelligi potest. Omnis
 autem essentia vel quiditas intelligi potest sine hoc quod alíquid intelligatur
 de esse suo facto; possum enim intelligere quid est homo ve! phoenix, ct tamen
 ignorare an esse habeant in natura. Ergo patet quod esse est aliud ab essentia
 vel quidditate» (De ente et Ess. c. V). — Como se vê, S. Tomás falla
 claramente da essência atual, e não da possível. Porquanto, ensina ele
 que «existência forma um composto com à essência; ora a existência, sendo
 coisa real, não pode formar o composto com uma coisa, que não existe na
 realidade, como é a essência possível.}
\pagmark{316}

     \medskip\noindent b) A existência distinguir-se-á realmente da essência, se
 . aquella não for um elemento constitutivo desta, e por isso se a
 essência existir em virtude de uma outra realidade
 ou perfeição;
 porque duas. coisas são realmente. distintas entre si, quando
 uma não é outra. Ora, no ente finito, a existência não é um
 elemento constitutivo da essência; porque, se o fosse; a criatura
 existiria em virtude da própria essência, isto é, seria essencial-
 mente existente, é por isso não seria contingente, mas necessária,
 e assim deixaria de ser criatura. Logo, no ente finito, a exi-
 stencia distingue-se realmente da essência \footnote{Este argumento pode propor-se por este outro modo:
 Se, no ente finito, a existência não se distinguisse realmente da essência,
 a existência das criaturas seria necessária, porque os elementos das essências
 são necessários; e por isso, assim como o homem é necessariamente racional
 (porque racional é um. elemento constitutivo da essência humana), a crea-
 tura seriá necessariamente existente. Ora a existência das criaturas não é
 necessária, mas contingente, porque podiam deixar de existir. Logo a existen-
 cia não é um elemento constitutivo da essência das criaturas, e por isso,
 no ente finito, a existência distingue-se. realmente da essência.
 Pode objetar se: Uma coisa, atualmente existente, sle necessária-
 mente, Logo não repugna que a essência atual das criaturas se diga e seja
 essencial ou necessariamente existente. Logo, no ente finito, a existência não
 se distingue realmente da essência. .
 Resposta. À necessidade na existência pode referir-se ao facto ou à exi-
 gencia. — Referida ao facto, à necessidade na existência é uma coisa hypo-
 thetica, e pode perfeitamente conciliar-se com a contingência do ente: de
 modo que um ente contingente, enquanto existe, não pode deixar de existir,
 existe necessariamente, e todavia não deixa de ser contingente, porque à exi-
 stencia não lhe era recessaria, isto é, porque não existe por assim o exigir a
 sua essência, — Referida à exigência, a necessidade na existência é uma coisa
 absoluta, constitui a própria essência dos entes, e por isso não pode conciliar-se
 com a contingência dos próprios entes: assim Deus, como veremos, existe por
 assim o exigir a sua essência, e por isso é o ser absolutamente necessário. —
 Ora a necessidade na existência, que se atribui aos entes finitos ou criados.
 não é uma necessidade da exigencia, mas é uma necessidade de facto; visto
 que os entes finitos existem, não porque o exige a sua essência, mas
 só porque Deus livremente os criou. Se os entes finitos existem, não porque
\pagmark{317}
 assio exige a sua essência, é claro gue a sua necessidade na existência não
é um elemento constitutivo da mesma essência. Se, nos entes Íinitos, x
necessidade da existência não é um elemento conslitutivo da essência, à existencia distingue-se realmente da essência. — Nem pode dizer-se que a necessidade ua existência se torna um elemento constilutivo da essência dos entes
finitos pelo facto destes existirem atualmente; porque, se assim fosse, as
essências não seriam-ijmmutaveis as-creaturasexistiriam-cm-virtude-da-sua..
essência e deixariam de ser contingentes: o que é absurdo. Logo a essência
dos seres finitos recebe ou tem a existência, mas não é a existência.. 
\pagmark{317}
  e) Todo o ato, que é recebido numa potência, distingue-se
realmente da potência, em que é recebido; porque repugna que
uma coisa seja recebida em si mesma. Ora, no ente finito, a
existência é recebida na essência, como o ato na potência; por
quanto, se a existência não fosse recebida na essência, como o
ato na potência, a existência dos entes finitos seria um ato
puro, subsistente, e por isso seria infinito e único: o que repugna. Logo, no ente. linito, a existência distingue-se realmente
da essência \footnote{Depois do que deixamos dito neste e no precedente artigo, não é
necessario insistir muito na verdade — que O ser (ou a existência) dos entes
finitos, se não fosse recebido na essência, como ato na potência, seria um
ato puro, subsistente, e por isso infinito e único. Porquanto, todo O ato É
limitado e multiplicado pela potência. que o recebe; e por isso um ato, que
não é recebido na potência, não pode ser limitado, nem multiplicado: mas é
necessariamente infinito e único. Oro O ser, ou a existência, dos entes finitos
é um ato essencialmente finito, multiplicável e multiplicado. — Nem se diga
que o ato, não recebido na essência, pode ser limitado ou multiplicado por
uma causa extrinseca; porque nenhuma coisa pode ser limitada ou multiplicada por uma cansa extrinseca, se é intrinsecamente ilimitada e única, e O
ato subsistente não tem em si mesmo um princípio, que o possa limitar ou
multiplicar. Diz o Angélico S. Tomás: «Esse subsistens oportet esse infinitum, quia mon terminatur aliquo recipiente... Impossibile est quod sit duplex
esse ommino infinitum, Esse enim. quod est infinitum, omnem perfectionem
essendi comprehendit ; et sic, si duobus talis adesset infinitas, non invenirelur quo num ab altero differret ((. Gent.. 1. IN, 52). — Este argumento
pode resumir-se assim: No ente finito, a existência refere-se à essência, como
o ato à potência. Ora o ato distingue-se realmente da potência. Logo, no
ente finito, a existência distingue-se realmente da essência.
 Em poucas palavras: Não podem realmente identificar-se duas coisas,
que são opostas entre si; aliás o ente e o não-ente seriam uma e a mesma
coisa na realidade. Ora a essência € O ser, nas criaturas, opõem-se por muitos modos, como se opõem o necessario e O contingente, O eterno e O temporuneo, O limitante € O limitado, O divisor e o diviso, O não-causado e o causado. Logo, nas criaturas, a essência e O ser distinguem-se realmente.
 Nas obras do S. Doutor encontram-se muitos outros argumentos, que
provam a distinção real entre a existência e a essência nos entes finitos,
\pagmark{318}
Cf. C. Gent., II, 52; Sum. Th., p. 1, q. 3, a. 4; q. 50, a. 2 ad 3; q. 75, a. 5 ad 4; q. 88, a. 2 ad 4; De Pot., q. 7, a. 2 ad 9; De Ver., q. 27, a. 1 ad 8; Com. in lib. De Hebd., l. 2). — Não insistimos; porque os argumentos
aduzidos são suflicientes para à demonstração da these.

Para a exacta compreensão da nossa these e para à refutação dus
dificuldades fazemos as seguintes considerações
 : .

     \medskip\noindent I. É certo que uma essência atual, por isso mesmo que é atual, existe.
Mas a questão consiste em saber se à atualidade compete a essa essência.
por si mesma, enquanto é essência, ou em virtude de uma outra realidade,
realmente distinta d" ella, que é a existência. Por quanto, uma coisa é e
diz-se atual ou porque é ato, ou porque recebe o ato, Ora a essência atual
— "não éachial-por
 siTmesma;-como-se
 ella fosse "ato, ou se identificasse
com O ato, — mas é atual.enquanto: recebe: 0 ato, que a completa na
linha da realidade e que é a existência. Sendo atual pelo facto de receber
O ato,
 a essência conserva sempre a razão ou a natureza de potência, e por
isso distingue-se realmente da existência, porque O ato e a potência são
coisas realmente dislinctas. — Se, às vezes, dizemos que a essência «atual
é atual por si mesma, comparamos a essência atual com a essência posssível, e não com a existência; porque, efetivamente, a essência atual
refere-se à essência possível, como O ato se refere à potência.

     \medskip\noindent II. A essência atual, ou física, é dotada de entidade própria, e tem
a sua última atualidade na própria linha, — na linha da essência, emquanto nada lhe falta do que é denotado pela sua definição. Mas esta última
atualidade na linha da essência não
 é suficiente para explicar a sua presença na realidade, Diz S. Tomás: « Natura considerata praeter suum esse,
est secundum acteptionem rationis tantum » (In | Sent. dist. 14, q.1, a. 3).
Se, pois, a essência existe na realidade, existe em virlude de uma outra
atualidade, — da atualidade de uma ordem diversa, que é o ato da existencia. — A existência, nos entes finitos, não é nem pode ser um elemento
constitutivo da essência; e por isso, quando se une com a essência, não a
altera, não a transforma, deixa-a ficar o que era; limita-se a completal-a na
ordem da existência, colocando-a e conservando-a-colocada na realidade.

     \medskip\noindent III. A existência da criatura, por isso mesmo que é existência, não é
indiferente para o opposto, isto é para à não-existência; aliás uma coisa
seria e deixaria de ser ao mesmo tempo e sob o mesmo aspecto. Por isso,
se uma coisa é indiferente para a existência e para a não-existência, esta
indiferença não deriva da própria existência, mas deriva do sujeito, em que
a existência é recebida e que pode existir e deixar de existir. — Como tambem, se a existência acaba, não acaba porque é existência, mas porque
acaba o seu sujeito. A razão é porque um ato misto, como é à existência
das criaturas, segue a sorte do sujeito, de que é ato.

     \medskip\noindent IV. O termo adequado da produção não é só à essência, nem é só a
\pagmark{319}
existência, mas é O composto de uma e de outra. Na verdade, a causa efi-
 ciente, quando produz a essência e à existência, as produz como duas
 coisas, mas como uma coisa só, composta de dois princípios, um dos quaes
 se refere ao outro, como a potência ao ato. — Por isso. nem a existência
 precede, na ordem do tempo, à essência, nem a essência precede a existência;
 mas, pelo facto mesmo-de-dar-a-existência-i-essenciazo-agente-não-só pri
 duz a existência mas também a essência. Diz S. Tomás: « Ex hoc ipso quod
 quidditati esse tribuitur, nou solum esse sed ipsa quiditas creari dicitur;
 quia antequam esse habeat, nibil est, nisi forte in intellectu creantis, ubi
 non est criatura, sed creatrix essentia
 » (De Pot., q. 3, a. 5 ad 2).

     \medskip\noindent V. D'onde se segue que a existência e a essência, embora sejam coisas
 realmente distintas, contudo não se podem separar uma de outra, sem que.
 uma e outra immediatamente acabem. Por quanto, o efeito, que de um
 ato deriva para um sujeito, não permanece e não pode permanecer, se
 acabar o próprio ato. Ora existir é O efeito próprio da existência, e esse
 efeito não pode ser produzido nem executado senão na essência. Repugna,
 pois, que a essência dure sem a existência; como repugna que a existência
 subsista sem a essência, para a qual tem uma relação necessária, — Nunca
 devemos esquecer que o ente existente é O composto, e que a existência c à
 essência não são duas realidades subsistentes, mas são os princípios intrin-
 sevos, — aclual e potencial, — pelos quaes O composto existe.. 
\pagmark{318}

\medskip\noindent\textbf{45. No Ente infinito, a existência não se distingue real-}
mente da essência.


     \medskip\noindent a) Uma coisa não se distingue realmente de outra, quando
entre elas ha perfeita identidade. Ora, no Ente infinito, ha perfeita identidade entre a existência e a essência; porque Deus
\pagmark{319}
existe em virtude ou pela exigencia da sua essência, e por isso
 a existência constitui a sua essência. Logo, no Ente infinito, a
 existência não se distingue realmente da essência.

     \medskip\noindent b) Se, no Ente infinito, a existência, se distingnisse real
 mente da essência, a existência seria recebida na essência, seria
 ato misto, não-subsistente, e por isso limitado, e multiplicado
 pela mesma essência, em que é recebido. Ora a existência, ou
 o ser de Deus é ato puríssimo, subsistente, e por isso infinito
 e único na sua ordem. Logo, no Ente infinito, à existência não
 se distingue realmente da essência \footnote{(1) Deus, e só Ele, é o ente necessario, que existe em virtude da sua
 essência, isto é, por assim O exigir a sua essência; e por isso em Deus, e só
 nEle, a existência não se distingue realmente da essência, mas identifica-se
 com esta. De facto, todos dizem — que Deus, e só Ele, é, — que Deus é o
 próprio Ser, a própria Existência,
 — que todos os seres finitos não são o ser,
 mas receberam O ser, ou à existência. .
 D" onde se vê que a essência pode realizar-se ou realizar-se por dois mo-
 dos: por si mesma, ou pela acção de uma causa extrinseca. Realiza-se por si
 mesmo, quando é necessária, quando existe na realidade por virtude própria,
 por assim ella o exigir; assim existe a Essência de Deus. Realiza-se pelo acção
 de uma causa extrinseca, quando é contingente, isto é. indiferente para a
 existência e para a não-existência, e por isso deve ser determinada por uma
"causa extrinseca; assim existe a essência do ente finito. — Quando à essen-
 cia se realiza por si mesma,  a existência é um ato constitutivo da mesma
 essência, isto é, a essência não só existe, mas é a sua existência, o seu ser.
\pagmark{320}
 Quando à essência é determinada para à existência pela acção de uma causa
extrinseca, à existência não é um elemento constitutivo da essência, mas é
"uma perfeição, que se lhe une, como O ato sé une com a potência, e que
lhe di o último complemento, a última atualidade na linha da realidade.
 Tal é a doutrina de S. Tomás. O S. Doutor não se cansa de repetir
 —
que em Deus O ser e à essência são uma € a mesma coisa, — que Deus é ato
subsistente, forma puríssima, não recebida em nenhuma potência, — e que
por isso Deus é infinito na sua perfeição e único na sua ordem, 30 passo, que,
vas criaturas, O ser é uma coisa diversa da essência, O ser é recebido na
essência, e por isso as criaturas são finitas e muitas. Eis as suas palavras: « Esse est atualitas ômnis formae vel naturae; non enim bonitas vel
humanitas significatur in actu, nisi prout significamus eam esse, Oportel
 igitur quod ipsum esse comparetur ad essentiam, quae est aliud ab ipso,
 sicut actus ad potentiam. Cum igitur in Deo nihil sit potentiale, sequitur
 quod non sit aliud in eo essentia quam suum esse. Sua igitur essentia est
suum esse. (Sum. Theol., p. 1, q.3, a. 4). — Mais adiante: « Infinitum, quod
 se tenet ex parte formae, non determinata per materiam, habet rationem
 perfecti.... Illud autem, quod est maxime formale omnium, est ipsum esse.
 Cum igitur esse divinum non sit esse receptum in aliquo, sed ipse sit suum
esse subsistens, manifestum est quod ipse Deus sit infinitus el perfectus
 »
 (Sum. Th. p. 1; q. 7, a. 1). — Noutro lugar diz: «Talis res, quae sit suum
esse, non potest esse nisi una; unde oportet quod in qualibet alia re, praeter
cam, sil aliud suum esse et aliud quiditas» (De Ente et Essentia, c. V). —
 « Hujusmodi formae (Angeli) non terminantur negue confrahuntur per aliquam materiam, Sed quia forma creata sic subsisteng habet esse, et nom est
suum esse, necesse est quod ipsum ejus esse sil receptum el contractum ad
 Lerminatam maturam. Unde non potest' esse infinitum simpliciter» (Sum.
 Th., p. 1, q. 7, a. 2). — « Esse hominis terminatum est ad hominis speciem,
quia est receptum in natura speciei humanae. Esse autem Dei, cum nom sit
in aliquo receptum, sed sit esse purum, non limitatur ad aliquem modum
perfectionis» (De pot., q. 1, a. 2: cf. De spirit. creat:, a. 1). — Nos citados
 trechos, que se referem aos entes finitos ou criados; trata-se de uma distinção real entre a essência e à existência, porque, se se tratasse apenas de
uma distinção lógica, não seria legitima a argumentação do S. Doutor:
visto que, também em Deus, ha uma distinção lógica entre a essência e à
existência. — Mas, deixando outras citações (Cf. 1 Sent., dist. 19, q.2, a. 2;
 C. Gentes, l. II, c. 52), O trecho seguinte torna a coisa evidente. No Comm.
ao livro de Boecio De Hebdomad. escreve: «Est ergo considerandum quod
sicut esse et quod est differunt in simplicibus secundum intentionem, ita in
compositis differunt realiter» (Lect. 2). O que se entende por ente simples,
dito com as seguintes palavras: «Hoc autem simplex unum et sublime est"
 ipse Deus», Poderia ser mais claro? Em Deus, que é o único ente simplíssimo, a essência é à existência diferem logicamente; em Lodas as crea-
\pagmark{321}
 turas, que são entes compostos, à essência e a existência diferem realmente (Cf, De Verit., q. 27, a. 1 ad 8).
 A distinção real, pois, entre a essência € à existência, nas criaturas, é
doutrina certamente propugnada pelo Angélico Mestre e é um dos pontos
capitais e fundamentais de toda a sua metafísica, como afirmam, entre
outros, o Card. Pallavicino (De Deo, c. IIl) e Silvestre Mauro (Quaest. phil.,
vol. II, q. 7). — Poderá impugnar-se esta doutrina; mas não pode, sem
grande temeridade, afirmar-se que é falso ou incerto que ella tenha sido
ensinada por 8. Tomás. Tiphano. que todavia não segue o angelico Mestre neste ponto. escreve: «Quod existentiam in creaturis re distinguat
(S. Tomás) ab essentia, id. vel negare, vel in dubium revocare, est hominis
impudentis, vel in ejus doctrina peregrini» (De hypostasi, c. 6).. 
\pagmark{320}

\medskip\noindent\textbf{46. Importancia da doutrina exposta. — À doutrina expo-}
 sta tem uma grande importancia, porque serve para estabelecer
uma harmonia perfeitíssima entre todos os seres. A simplicidade
é signal de perfeição, como a composição é signal de limitação
ou imperfeição. Deus é ato simplíssimo, puríssimo, a sua exi-
\pagmark{321}
 stencia identifica-se com a essência; por isso é infinito na perfeição.
— As inteligências separadas, que são os Anjos, são compostas
de essência e de existência; por isso são limitadas na perfeição,
infinitamente inferiores a Deus, nas superiores a todas as outras
criaturas. — As criaturas visiveis são compostas de essência e
de existência, e também de matéria e de forma; por isso são
ainda mais limitadas, e ínfimo é o seu grau de perfeição \footnote{(1) Nem essa é a única vantagem. A questão acerca da distinção real
entre a existência e à essência, nas criaturas, liga-se intimamente com as
outras duas importantes questões acerca de distineção real entre q substância
e os acidentes, e entre a materia e a forma. Estas tres questões reflectem,
pela sua vez, muita luz sobre alguns dogmas da Fé. Na verdade, a teoria
da distinção real entre a essência e à existência explica, de algum modo,
como é que, em N. S. Jesus Cristo, podem duas: naturezas ou essências
subsistir numa só existência, ou pessoa divina: — à teoria da distinção
real entre a substância e os acidentes serve para conhecer, de algum modo,
o mistério da Eucaristia: — a teoria da materia e da forma explica a
união da alma com o corpo em harmonia com as definições da Igreja.}

\end{document}
